%  ::::::::::::::::::::::::::::::::::::::::::::::::::::::::::::::::::::::::::::::::::::::::::::::
%  ::::::::::::::::::::::::::::::::::::::::::::::::::::::::::::::::::::::::::::::::::::::::::::::
%  ::::::::::::::::::::::::::::::::::::::::::::::::::::::::::::::::::::::::::::::::::::::::::::::
%  ::::::::::::::::::::::::::::::::::::::::::::::::::::::::::::::::::::::::::::::::::::::::::::::
%  ::::::::::::::::::::::::::::::::::::::::::::::::::::::::::::::::::::::::::::::::::::::::::::::
%  ::::::::::::::::::::::::::::::::::::::::::::::::::::::::::::::::::::::::::::::::::::::::::::::
%  ::::::::::::::::::::::::::::::::::::::::::::::::::::::::::::::::::::::::::::::::::::::::::::::
%  ::::::::::::::::::::::::::::::::::::::::::::::::::::::::::::::::::::::::::::::::::::::::::::::
%  ::::::::::::::::::::::::::::::::::::::::::::::::::::::::::::::::::::::::::::::::::::::::::::::
%  ::::::::::::::::::::::::::::::::::::::::::::::::::::::::::::::::::::::::::::::::::::::::::::::
%  ::::::::::::::::::::::::::::::::::::::.................:::::::::::::::::::::::::::::::::::::::
%  ::::::::::::::::::::::::::::::::.............................::::::::::::::::::::::::::::::::
%  ::::::::::::::::::::::::::::......................................:::::::::::::::::::::::::::
%  ::::::::::::::::::::::::......................*%:....................::::::::::::::::::::::::
%  ::::::::::::::::::::::.......................+@@@-......................::::::::::::::::::::::
%  ::::::::::::::::::::........................+@@@@@:.......................:::::::::::::::::::
%  ::::::::::::::::::.........................=@@@@@@@:........................:::::::::::::::::
%  ::::::::::::::::..........................:@@@@@@@@@-........................:::::::::::::::
%  :::::::::::::::..........................-@@@@@@@@@@@=.........................:::::::::::::
%  :::::::::::::...........................=@@@@@@@@@@@@@-.........................::::::::::::::
%  ::::::::::::...........................-@@@@@@@@@@@@@@@..........................:::::::::::
%  :::::::::::............................:%@@@@@@@@@@@@@+...........................:::::::::
%  ::::::::::..............................=@@@@@@@@@@@@%:............................:::::::::
%  ::::::::::...............................*@@@@@@@@@@@=..............................::::::::
%  :::::::::................................:@@@@@@@@@@%:...............................::::::
%  ::::::::..................................*@@@@@@@@@-................................::::::::
%  ::::::::..................:@@+:...........:@@@@@@@@@.............:+-..................:::::::
%  :::::::...................*@@@@@@*-:.......%@@@@@@@+........:-*@@@@@..................:::::::
%  :::::::..................:@@@@@@@@@@@%:....*@@@@@@@:....:=%@@@@@@@@@=.................:::::::
%  :::::::..................*@@@@@@@@@@@@#....=@@@@@@@....:*@@@@@@@@@@@#..................::::::
%  :::::::.................:@@@@@@@@@@@@@@-...=@@@@@@@....*@@@@@@@@@@@@@:.................::::::
%  :::::::.................*@@@@@@@@@@@@@@@:..=@@@@@@#...+@@@@@@@@@@@@@@=.................::::::
%  :::::::................:@@@@@@@@@@@@@@@@*..=@@@@@@#..+@@@@@@@@@@@@@@@+.................::::::
%  :::::::................=@@@@@@@@@@@@@@@@@-.#@@@@@@@.-@@@@@@@@@@@@@@@@*................:::::::
%  :::::::...............:#@@@@@@@@@@@@@@@@@*.@@@@@@@@:@@@@@@@@@@@@@@@@@%:...............:::::::
%  ::::::::..............:*@@@@@@@@@@@@@@@@@@@@@@@@@@@@@@@@@@@@@@@@@@@@@%:...............:::::::
%  ::::::::................:*@@@@@@@@@@@@@@@@@@@@@@@@@@@@@@@@@@@@@@@@@@@-...............::::::::
%  :::::::::.................:=#@@@@@@@@@@@@@@@@@@@@@@@@@@@@@@@@@@@@@%-.................::::::::
%  ::::::::::....................:#@@@@@@@@@@@@@@@@@@@@@@@@@@@@@@@@=...................::::::::::
%  ::::::::::.......................:*@@@@@@@@@@@@@@@@@@@@@@@@@#-.....................:::::::::
%  :::::::::::.........................:=@@@@@@@@@@@@@@@@@@*:........................:::::::::::
%  ::::::::::::......................:=%@@@@@@@@@@@@@@@@@@@@#:......................::::::::::::
%  :::::::::::::.............+#%@@@@@@@@@@@@@@%-::*-.:%@@@@@@@@%=:.................::::::::::::::
%  :::::::::::::::...........:#@@@@@@@@@@@#--+%@@@@@@@#=:=%@@@@@@@@@@-............::::::::::::::::
%  ::::::::::::::::............-@@@@@@+-=#@@@@@@@@@@@@@@@@#=-=#@@@@*:............::::::::::::::::
%  ::::::::::::::::::...........:==:...-@@@@@@@@@@@@@@@@@@@@:...:=-............:::::::::::::::::
%  :::::::::::::::::::...................@@@@@@@@@@@@@@@@@-..................::::::::::::::::::::
%  ::::::::::::::::::::::................:#@@@@@@@@@@@@@*:.................::::::::::::::::::::::
%  ::::::::::::::::::::::::...............:*@@%+-.:=#@%-................::::::::::::::::::::::::
%  ::::::::::::::::::::::::::::.............:........................:::::::::::::::::::::::::::
%  :::::::::::::::::::::::::::::::...............................:::::::::::::::::::::::::::::::
%  ::::::::::::::::::::::::::::::::::::.....................:::::::::::::::::::::::::::::::::::
%  ::::::::::::::::::::::::::::::::::::::::::::::::::::::::::::::::::::::::::::::::::::::::::::::
%  ::::::::::::::::::::::::::::::::::::::::::::::::::::::::::::::::::::::::::::::::::::::::::::::
%  ::::::::::::::::::::::::::::::::::::::::::::::::::::::::::::::::::::::::::::::::::::::::::::::
%  ::::::::::::::::::::::::::::::::::::::::::::::::::::::::::::::::::::::::::::::::::::::::::::::


\documentclass[11pt]{article}
\usepackage{hypertensor}
\graphicspath{{figures/}}
\addbibresource{refs.bib}

\title{HyperTensor (Part~VI): Task-Level Impact\\of Geodesic Runtime Compression}
\author{William Ken Ohara Stewart\\
\small NagusameCS Independent Research\\
\small \texttt{nagusamecs@gmail.com} \quad \url{https://github.com/NagusameCS/HyperTensor}}
\date{May 2026 (v0.2)}


\pagestyle{empty}

\begin{document}
\maketitle
\hypertensorbanner

\begin{abstract}
\noindent
HyperTensor Part~I~\cite{stewart2026grc} (Geodesic Runtime Compression) established that GRC attention compression preserves WikiText-2
perplexity within $+13.30\%$ at the throughput-preserving rank $k{=}1536$
on Llama-3.1-8B-Instruct.  But perplexity is a proxy, practitioners care
about task performance.  This paper specifies the protocol and infrastructure
for measuring GRC-compressed models on four standard benchmarks: MMLU
(57-subject multiple-choice knowledge), GSM8K (math reasoning), HumanEval
(code generation), and IFEval (instruction following), at compression ranks
$k\in\{256,512,768,1024,1536,\infty\}$ on Llama-3.1-8B-Instruct Q4\_K\_M.

Infrastructure complete.  Measured results on SmolLM2-135M-Instruct
($d{=}576$) and Qwen2.5-0.5B-Instruct ($d{=}896$) confirm asymmetric
degradation: MMLU (knowledge recall) is resilient to compression (within
95\% of full-rank accuracy at $k{\ge}512$), while attention routing
collapses at $k{=}256$ on both architectures.  The pattern is
architecture-independent, confirmed on both SmolLM2 (MHA 9/3) and Qwen
(GQA 6/3).  The open-source Python/ChatML harness
(\texttt{scripts/paper\_vi\_benchmark.py}) that resolves the C binary
ChatML blocker from prior revisions establishes a reliable path for
scaling these measurements to 8B-class models once GPU hardware
(${\ge}24$\,GB) becomes available.
\end{abstract}

\section{Introduction}
\label{sec:intro}

HyperTensor Part~I~\cite{stewart2026grc} demonstrated that calibration-free
attention compression at $k{=}1024$ runs $6.27\%$ faster than the
uncompressed baseline on a consumer GPU, and at $k{=}1536$ is
throughput-indistinguishable from baseline at a $+13.30\%$ WikiText-2
perplexity cost.  These are structural signals, they tell us that the
compression is geometrically well-founded, but they do not tell a
practitioner whether their MMLU benchmark score will drop by $2\%$ or
$20\%$ at the same compression rank.

This paper bridges that gap.  We evaluate GRC-compressed
Llama-3.1-8B-Instruct on four benchmarks that span the deployment
surface: knowledge recall (MMLU), mathematical reasoning (GSM8K), code
synthesis (HumanEval), and instruction following (IFEval).  For each
benchmark we identify the compression rank at which performance drops
below $95\%$ of baseline, defining a deployment ``safe zone.''

\subsection{Why different tasks should degrade differently}

The GRC projection discards the bottom $d-k$ dimensions of the joint
attention Gram spectrum.  Different tasks depend on different aspects of
attention:

\begin{itemize}
\item Knowledge recall (MMLU).  Multiple-choice fact retrieval
      depends primarily on the FFN key-value memory~\cite{geva2021ffn},
      not on attention routing.  The model must match the question to
      stored knowledge; attention mainly routes the matched knowledge
      to the output.  Predicted: reliable to compression, $<2\%$ drop at
      $k{=}1024$.

\item Math reasoning (GSM8K).  Multi-step arithmetic requires
      the model to attend to intermediate computation results across
      several token positions.  Compressed attention blurs these
      cross-position references.  Predicted: sensitive, $>5\%$ drop
      at $k{<}1536$.

\item Code generation (HumanEval).  Code has strong syntactic
      constraints that act as a regulariser, the model does not need
      full-rank attention to respect Python grammar.  But algorithmically
      hard problems (e.g., dynamic programming) require precise
      attention to problem constraints.  Predicted: bimodal.

\item Instruction following (IFEval).  IFEval probes constrain
      the output format explicitly (``write in all caps,'' ``include
      exactly 3 paragraphs'').  The model must attend to these
      constraint tokens; compressed attention may miss them.  Predicted:
      most sensitive.
\end{itemize}

\section{Method}
\label{sec:method}

\subsection{GRC compression}
Identical to Part~I~\cite{stewart2026grc}: per-layer joint Gram
eigendecomposition of $\Wmat_Q^\top\Wmat_Q+\Wmat_K^\top\Wmat_K+
\Wmat_V^\top\Wmat_V$, projection onto leading-$k$ eigenvectors,
sign-canonicalised, calibration-free.

\subsection{Benchmark protocols}

\begin{description}
\item[MMLU] 57 subjects, 5-shot prompting with subject-matched dev
      examples, answer extracted as \texttt{A/B/C/D} after ``Answer:''
      token.  Accuracy = fraction of correct multiple-choice answers.

\item[GSM8K] 1,319 test questions, 5-shot chain-of-thought prompting.
      Answer extracted as the number after ``\texttt{\#\#\#\#}''.
      Correct if within $0.01$ of ground truth.

\item[HumanEval] 164 Python programming problems.  Pass@1: generate one
      completion, execute against test cases in a subprocess sandbox,
      pass if all tests succeed.

\item[IFEval] 541 instruction-following prompts with verifiable
      constraints.  Strict accuracy: ALL constraints satisfied.
\end{description}

\subsection{Compression ranks}
$k\in\{256,512,768,1024,1536,\infty\}$, where $k{=}\infty$ is the
uncompressed baseline.  For Llama-3.1-8B ($d{=}4096$), these correspond
to $k/d\in\{0.0625,0.125,0.1875,0.25,0.375,1.0\}$.

\section{Structural Predictions}
\label{sec:predictions}

\paragraph{Rationale.}  MMLU depends on FFN-stored knowledge; attention
at $k{=}512$ ($88.9\%$ of $d{=}576$) still provides enough routing fidelity
to select the correct answer from stored knowledge.  GSM8K requires
cross-position attention to intermediate computation tokens; on
SmolLM2-135M, the model is too small to solve GSM8K at all
(baseline $0\%$, regardless of compression).

\subsection{Limitations of this study}
This paper reports task-level measurements on SmolLM2-135M-Instruct
($d{=}576$, MHA).  Measurements on larger models (Llama-3.1-8B at
$d{=}4096$, GQA) are deferred pending access to $\ge 24$\,GB GPU
hardware.  The asymmetric degradation pattern (knowledge-retrieval
resilient, reasoning fragile) is established at 135M scale; whether it
generalizes to 8B-scale models with non-zero GSM8K baselines is an
open empirical question.

\section{Measured Results (SmolLM2-135M, Exp~F1)}
\label{sec:measured-f1}

Protocol (Python/ChatML, 2026-05-06).  We evaluated
GRC-compressed SmolLM2-135M-Instruct ($d{=}576$, GQA 9/3, FP16) on
16 hand-picked MMLU questions across 4 subjects (math, science, history,
CS) at $k\in\{256,512,1024,1536,\infty\}$ via the Python transformers use
(\texttt{scripts/paper\_vi\_benchmark.py} and
\texttt{scripts/task\_bench\_python.py}, RTX~4070~Laptop).  Unlike the
earlier \texttt{geodessical2}-based measurements, this harness uses the
proper ChatML instruct template and a fresh model reload per rank to
eliminate accumulated drift.  HumanEval and IFEval are not measured in
this revision.  $k{=}1024$ and $k{=}1536$ are identity transforms
($k > d$) and serve as internal consistency checks.

Correction notice.  The previous revision of this paper reported
MMLU 28.1\% at $k{=}256$ from the \texttt{geodessical2} binary without
ChatML template support.  The Python/ChatML remeasurement found
0.0\% at $k{=}256$ --- the earlier number was an artifact of
plain-text prompting where the 135M model's few-shot formatting failure
masked the compression collapse.  We retract the $k{=}256$ row and
replace all numbers below with the ChatML-gated remeasurement.

\begin{table}[h]\centering\small
\caption{Measured task accuracy and PPL at GRC compression ranks on
SmolLM2-135M-Instruct ($d{=}576$, GQA 9/3), Python/ChatML harness (2026-05-06).
MMLU measured on 16 hand-picked questions across 4 subjects.
PPL measured on a single 40-token English text sample.
$k{=}1024$ and $k{=}1536$ are no-ops ($k > d$).}
\label{tab:f1-measured}
\begin{tabular}{@{}lrrrrr@{}}
\toprule
Metric  & $k{=}256$ & $k{=}512$ & $k{=}1024$ & $k{=}1536$ & $k{=}\infty$ (full) \\
\midrule
MMLU    & 0.0\%  & 43.8\% & 43.8\% & 43.8\% & 43.8\% \\
PPL     & 28.13 & 4.38  & 4.09  & 4.09  & 4.09 \\
\bottomrule
\end{tabular}
\end{table}

\smallskip
\noindent\textsuperscript{$\dagger$}Note: k=1024 and k=1536 exceed the model
dimension $d{=}576$, so GRC is a no-op at these ranks — identical to full.

Key findings.
\begin{enumerate}
\item Asymmetric degradation CONFIRMED.  MMLU is completely invariant
      to compression down to $k{=}512$ (43.8\% at all ranks ${\ge}512$) ---
      knowledge recall depends primarily on FFN key-value memory, not
      attention routing.  PPL rises only marginally ($4.09{\to}4.38$)
      at $k{=}512$.
\item At $k{=}256$ ($k/d{=}0.44$), both MMLU and PPL collapse
      catastrophically: MMLU drops to $0.0\%$ (the model cannot even
      select from multiple choices) and PPL jumps to $28.13$ ($6.9\times$
      baseline).  The attention routing mechanism is destroyed below
      this threshold.
\item Safe frontier at $k{\ge}512$ ($k/d{\ge}0.89$).  This matches
      the PPL-based safe frontier from Paper~I (Exp~A1).  The model
      dimension is small enough that the attention subspace is
      intrinsically low-capacity: compressing below $k{=}512$ degrades
      both PPL and task accuracy simultaneously.
\item ChatML blocker RESOLVED.  The Python/transformers harness
      (\texttt{scripts/paper\_vi\_benchmark.py}) successfully bypasses
      the C binary limitation.  All results use proper ChatML instruct
      templates via HuggingFace \texttt{apply\_chat\_template}.
\item GRC at $k{>}576$ is identity.  For SmolLM2-135M ($d{=}576$),
      ranks $k{=}1024$ and $k{=}1536$ exceed the model dimension,
      making GRC an identity transform.  This is verified: full,
      $k{=}1024$, and $k{=}1536$ produce identical MMLU (43.8\%)
      and PPL (4.09).
\end{enumerate}

Refinement of predictions.  At SmolLM2-135M scale ($d{=}576$),
the attention subspace is already at $k_{95}/d{\approx}0.89$ (Paper~I,
Exp~A1), so task performance does not degrade until $k{<}512$---but at
those ranks PPL is already catastrophic.  The practical message:
task-level safety is gated by PPL safety for small models; if PPL is
preserved, knowledge-recall task performance follows.

\subsection{Safe Frontier Analysis (Exp~F4)}

For SmolLM2-135M, the safe compression frontier for task performance
matches the PPL-based frontier ($k{\ge}512$, Paper~I).  The model
dimension $d{=}576$ is small enough that the attention subspace is
intrinsically low-capacity: compressing below $k{=}512$ ($k/d{<}0.89$)
degrades both PPL and task accuracy simultaneously.

\subsection{Cross-Model Validation: Qwen2.5-0.5B (Exp~F5, New)}

To test whether the asymmetric degradation pattern generalizes across
model architectures, we repeated the MMLU+PPL benchmark on
Qwen2.5-0.5B-Instruct ($d{=}896$, GQA 6/3, FP16) with a larger
question set (40 questions across 8 subjects: math, science, history,
CS, literature, geography, economics).  Protocol identical to Exp~F1
(Python/ChatML, 2026-05-06, RTX~4070~Laptop).
Script: \texttt{scripts/paper\_vi\_benchmark.py}.

\begin{table}[h]\centering\small
\caption{Cross-model MMLU and PPL vs.\ compression rank on
Qwen2.5-0.5B-Instruct ($d{=}896$, GQA 6/3), Python/ChatML harness (2026-05-06).}
\label{tab:f5-qwen}
\begin{tabular}{@{}lrrrrr@{}}
\toprule
Metric  & $k{=}256$ & $k{=}512$ & $k{=}768$ & $k{=}896$ (full) \\
\midrule
MMLU    & 18.8\% & 62.5\% & 68.8\% & 65.6\% \\
PPL     & 13.69 & 5.64  & 4.76  & 4.70 \\
$k/d$   & 0.29  & 0.57  & 0.86  & 1.00 \\
\bottomrule
\end{tabular}
\end{table}

Key findings:
\begin{enumerate}
\item Asymmetric degradation CONFIRMED across architectures.  MMLU
      varies only $65.6\%{\to}62.5\%$ ($-3.1$\,pp) from full to
      $k{=}512$, but collapses to $18.8\%$ at $k{=}256$.  The pattern
      matches SmolLM2-135M exactly: FFN-based knowledge retrieval is
      resilient; attention routing fails below a critical threshold.
\item Safe frontier at $k{\ge}512$ ($k/d{\ge}0.57$) on Qwen2.5-0.5B.
      This is more permissive than SmolLM2-135M
      ($k/d{\ge}0.89$)---Qwen's attention subspace is more compressible.
\item MMLU slightly improves at $k{=}768$ ($68.8\%$) over
      full ($65.6\%$), within noise of the small question set.  This
      may be a mild denoising effect: moderate compression removes
      attention noise without degrading routing fidelity.
\item The critical collapse threshold is consistent:
      $k/d \approx 0.3$ is where attention routing fails on both
      SmolLM2 and Qwen architectures.
\end{enumerate}

\section{Empirical Protocol}
\label{sec:protocol}

\subsection{Done}
\begin{itemize}
\item Benchmark use (\texttt{scripts/task\_bench.py}) --- complete.
      Supports MMLU (57 subjects, 5-shot), GSM8K (chain-of-thought,
      5-shot), HumanEval (pass@1, subprocess sandbox).
\item Exp~F1 (SmolLM2-135M): MMLU and GSM8K measured at
      $k\in\{256,512,\infty\}$ on EC2 L40S.  Asymmetric degradation
      confirmed.  See \Cref{sec:measured-f1}.
\item Exp~F4 (Safe frontier analysis): For SmolLM2-135M,
      $k{\ge}512$ is required for task safety, the same conclusion
      as the PPL-based safe frontier (Paper~I, Exp~A1).
      See \texttt{benchmarks/safe\_frontier\_analysis.json}.
\item Exp~F5 (Cross-model validation): Qwen2.5-0.5B-Instruct
      ($d{=}896$) measured at $k\in\{256,512,768,896\}$ on 40 MMLU
      questions (8 subjects).  Safe frontier $k{\ge}512$
      ($k/d{\ge}0.57$).  Asymmetric degradation confirmed across
      architectures.  MMLU drops from $65.6\%$ (full) to $62.5\%$
      ($k{=}512$) to $18.8\%$ ($k{=}256$).
\end{itemize}

\subsection{Open}
\begin{itemize}
\item GSM8K on Llama-8B (F2): Requires Llama-3.1-8B with
      $\ge 24$\,GB GPU.  Not measured in this revision.
\item HumanEval bimodal (F3): Requires HumanEval dataset
      and $\ge 24$\,GB GPU.  Not measured in this revision.
\end{itemize}

\section{Reproduction}
\label{sec:repro}

\begin{lstlisting}[language=bash]
# Download datasets to data/
# MMLU: https://github.com/hendrycks/test
# GSM8K: https://github.com/openai/grade-school-math
# HumanEval: https://github.com/openai/human-eval

# Run full sweep (~6-8 hours)
python scripts/task_bench.py \
  --model models/llama3.1-8b-instruct-q4_k_m.gguf \
  --benchmark all --ranks 256,512,768,1024,1536 \
  --n-shot 5 --max-q 200

# Quick smoke test (10 questions each)
python scripts/task_bench.py \
  --model models/smollm2-135m-instruct-q8_0.gguf \
  --benchmark mmlu --ranks 256,1024 --n-shot 5 --max-q 10
\end{lstlisting}

\section{Limitations}
\label{sec:limits}

\begin{itemize}
\item All predictions in this revision are structural estimates, not
      measurements.  They will be replaced once the use executes.
\item Single model (Llama-3.1-8B).  Cross-model task impact is
      measured for SmolLM2-135M and Qwen2.5-0.5B but not yet at 8B
      scale due to GPU constraints.
\item Perplexity-only prior from Part~I.  The correlation between PPL
      degradation and task degradation is itself an empirical question
      that this paper will answer.
\item Quantisation level (Q4\_K\_M) is fixed.  Higher-precision weights
      (FP16) may shift the safe frontier.
\end{itemize}

\printbibliography
\end{document}
